\chapter{Review of Related Literature}
\label{ch:literature}

Activated sludge surrogate modeling draws on several bodies of research. Mechanistic models define the treatment processes and component balances to be approximated. Data-driven studies show how treatment quantities can be estimated from measurements or simulations. Physical enforcement methods restrict predictions to specified balances and concentration bounds. Interpretable modeling explains the fitted relationships, while optimization studies examine their use in operating decisions. This chapter reviews the contributions of these studies and the questions that remain when their methods are combined for activated sludge analysis.

\section{Mechanistic foundations of activated sludge modeling}

\subsection{Development of component and nutrient models}

Activated sludge modeling developed from the need to relate operating conditions to organic matter removal, oxygen demand, and sludge production. \citet{Marais1976} formulated steady-state activated sludge analysis around substrate utilization and biomass. \citet{Dold1981} developed a general model that linked several biological processes within a common treatment calculation. These contributions established a basis for calculating treatment responses from material balances and biological rates.

\citet{Henze1987} presented the single-sludge framework associated with Activated Sludge Model No.\ 1 (ASM1). The model distinguishes substrate, biomass, and nitrogen components so that carbon oxidation, nitrification, and denitrification can be calculated together. Later models addressed processes that this framework did not include explicitly. \citet{Henze1999} extended phosphorus-removal modeling to include the denitrifying activity of phosphorus-accumulating organisms. \citet{Gujer1999} introduced Activated Sludge Model No.\ 3 (ASM3), which includes organic substrate storage and uses endogenous respiration to model biomass decay.

These developments addressed different biological modeling needs. The phosphorus-removal extension connects nitrate use to phosphate behavior. The storage-based model separates substrate uptake from its later use by biomass. The collected formulations in \citet{Henze2006} provide a common reference for these component definitions and their associated reaction terms. Their contribution to surrogate modeling is a physically defined prediction target. Effluent COD or TN can summarize treatment quality, while the component concentrations describe the material forms underlying those totals.

\citet{Hauduc2013} compared seven activated sludge models by examining their treatment of hydrolysis, fermentation, microbial growth and decay, and biological storage. Their review identifies conceptual differences between the models and relates those differences to the underlying process knowledge. It shows that model selection depends on the treatment process and the question being studied. For a surrogate comparison, this means that all candidate predictors need to approximate the same selected process model. Otherwise, differences in prediction can arise from different biological assumptions as well as different statistical methods.

\subsection{Biological dependence and calibration uncertainty}

Nutrient-removal research explains why operating inputs and influent composition must be considered together. \citet{Wentzel1992} examined the interactions among nitrification, denitrification, and excess phosphorus removal. \citet{Guerrero2011} studied competition for organic substrate in a pilot system intended for simultaneous nitrogen and phosphorus removal. Their experiments found that the nature of the supplied carbon source affected the ability to maintain phosphorus removal when nitrate entered the initially anaerobic zone. The study therefore addresses a specific biological interaction between substrate composition, recycle conditions, and nutrient removal.

For statistical modeling, these findings support retaining both operating conditions and component-resolved influent information. A model based only on total organic loading can omit differences between carbon sources that have different biological roles. A model that treats aeration or recycle independently of loading can also overlook the conditions under which those actions affect nutrient removal. This is a reason to examine interactions in a treatment surrogate. It is not evidence that a fitted interaction coefficient identifies a microbial mechanism by itself.

Calibration research addresses another part of the modeling problem. \citet{Petersen2003} reviewed experimental designs for estimating activated sludge parameters. \citet{Vanrolleghem1999} examined the use of respirometry, which measures oxygen consumption, to estimate model parameters and component combinations. These methods help connect the internal quantities of a process model to observations. They also show why the information supplied by an experiment must match the parameter being estimated.

\citet{Machado2014} calibrated a phosphorus-removal model using data from a full-scale plant in Manresa, Spain. They assessed parameter confidence intervals and examined how influent and operating variables affected the model fit. They also varied those inputs while retaining default kinetic and stoichiometric parameters to evaluate output sensitivity. The study addresses the influence of input uncertainty alongside uncertainty in biological parameters. Its relevance to surrogate modeling is that simulated targets depend on both sets of assumptions. Agreement with those targets assesses the surrogate's approximation, while agreement with an operating plant requires calibration and observational evidence.

The mechanistic literature thus establishes meaningful components, reaction relationships, and methods for calibration. The unresolved question for this dissertation is how to approximate those defined treatment responses efficiently while retaining their material accounting. The surrogate is evaluated against an adopted mechanistic reference. It does not replace the separate task of calibrating that reference to a particular plant.

\subsection{Clarification and connected plant balances}

Biological reactors do not determine the final discharge independently of solids separation. \citet{Takacs1991} developed a clarification and thickening model based on solids fluxes and balances around layers of a one-dimensional settler. The model calculates the solids profile and the overflow and underflow concentrations under steady and changing conditions. \citet{HartelPopel1992} also modeled secondary clarification with sludge thickening. These studies provide mechanistic descriptions of how suspended material is retained, concentrated, and discharged.

\citet{JeppssonDiehl1996} examined the propagation of biological components through the secondary clarifier. They identified difficulties in coupling a biological reactor model to a settler model because the two models describe particulate material differently. Their evaluation of component-transport algorithms focused particularly on the composition of sludge returned to the reactor. This study directly establishes why predicting aggregate solids alone is insufficient for every connected component calculation.

\citet{Alex2008} addressed the need for comparable plant assessments through Benchmark Simulation Model No.\ 1. The benchmark connects five biological compartments, a secondary settler, internal recycle, return sludge, and wasting. It also specifies influent conditions, control settings, and assessment criteria. Its contribution is a common process and evaluation setting in which alternative operating strategies can be compared. It makes the connection between biological treatment and separation explicit.

Together, these studies establish the plant relationships that a connected surrogate must approximate. A clarifier prediction affects both the discharge and the recycled biomass entering the biological tanks. Stored solids also contribute to the plant inventory used in solids-retention calculations. A single-reactor surrogate therefore addresses only part of the plant response. Extending physical enforcement to a connected plant requires consistent component flows at shared streams and an explicit account of retained solids.

\section{Data-driven prediction of wastewater treatment responses}

\subsection{Soft sensors and reported treatment quantities}

Data-driven treatment studies developed partly because plants collect measurements that are difficult to interpret jointly. \citet{Newhart2019} reviewed methods for fault detection, variable prediction, and advanced process control. They identify missing observations, outliers, correlated variables, and changing process behavior as important features of plant data. \citet{Sundui2021} reviewed machine learning applications across biological wastewater treatment. These reviews establish several useful prediction tasks and explain the need to match a method to its data and intended application.

A soft sensor estimates a quantity that is difficult or slow to measure from other available observations. \citet{Ching2022} developed a soft sensor for five-day biochemical oxygen demand using extreme gradient boosting (XGBoost). They evaluated influent and effluent estimates at two treatment plants and used supporting measurements selected from variables such as COD, solids, nutrients, and simpler sensor readings. The model estimated a wide range of oxygen-demand values and could accommodate missing supporting measurements. The study addresses the practical need for faster monitoring of a selected water quality quantity.

\citet{Wang2025} compared a mechanistic activated sludge model with six machine learning models for effluent TN prediction using one year of high-resolution municipal plant data. Their results showed that the mechanistic model captured the direction of TN variation, while the learned models achieved greater predictive accuracy in the assessed data. They also used feature-attribution analysis to identify influential inputs. This provides evidence that statistical modeling and model inspection can improve a defined effluent-prediction task.

\citet{Ekinci2023} addressed waste-sludge prediction using 208 daily observations from an advanced biological treatment plant in Turkey. Their inputs included flow, polymer dosage, and the amounts of solids and pollutants removed. They compared learning methods and reported improved prediction after feature selection. The study adds a useful solids-management target to the quality indicators considered in other work. Its measured removal inputs support that estimation task. A prediction made before evaluating a proposed operating setting would require a different set of available inputs.

The common achievement of these studies is useful estimation of selected treatment quantities from available plant information. Their results support monitoring, diagnosis, and particular operating analyses. The remaining question here concerns a different output requirement. Predicting all modeled effluent components is needed when the analysis must check mass conservation, non-negativity, and the distribution of material among biological forms. Accuracy for TN, oxygen demand, or sludge production supplies part of that evidence, while complete component prediction requires its own assessment.

\subsection{Component estimation and simulation-based surrogates}

Component-estimation studies move closer to the information required by mechanistic models. \citet{Wang2024} used a backpropagation neural network to estimate five COD-related component quantities from measured influent COD. They first characterized the components in an operating plant and then evaluated the neural predictions against those measurements. Their reported prediction accuracy exceeded 80\% for the five assessed components. The study addresses the difficulty of repeatedly measuring component allocations needed as model inputs.

This contribution differs from predicting the effluent component state at a proposed operating setting. Influent fractionation estimates how the incoming organic material is divided among model categories. An effluent surrogate estimates the result of biological conversion under specified controls. The two tasks can complement one another, but they have different reference outputs and assessment requirements. The component estimates of \citet{Wang2024} therefore support improved influent characterization without establishing the complete treatment mapping investigated here.

Simulation-based studies address that mapping more directly. \citet{Fang2011} combined an extended activated sludge model with a support vector surrogate and a genetic optimization search for a full-scale nutrient-removal system. Mechanistic simulations supplied the relationship between operating conditions and effluent quality learned by the surrogate. Their optimization suggested that anoxic volume and internal recycle could be reduced while retaining the specified discharge quality. The study is an early example of combining process simulation, statistical approximation, and operating search.

\citet{He2023} developed a Kriging surrogate for greenhouse-gas emissions from papermaking wastewater treatment. Kriging uses statistical relationships among sampled input conditions to estimate responses at other inputs. They used steady-state process simulations to calculate emissions and fitted the surrogate to those outputs. The reported coefficient of determination exceeded 0.96, and the model supported analysis and optimization of controllable treatment variables. Their work addresses the cost of repeated evaluation of an environmental objective. Its learned target is the emissions measure, rather than the complete component response needed for the mass conservation checks of this dissertation.

\citet{Selerio2026} addressed complete single-reactor component prediction through invariant-constrained second-order regression. The study used 10,000 simulated steady-state cases with 20 effluent components and compared the structured regression with nine statistical predictors. Its evaluation included training-size variation, reported water quality errors, and physical violation diagnostics. This provides a closer precedent for the present component-level task than the selected-indicator studies. Its physical enforcement and interpretation are examined in the following sections.

Table~\ref{tab:rrl_prediction_evidence} summarizes the progression from monitoring to process approximation. The distinctions concern the purpose of the studies and the information their assessments establish. They identify the additional evidence required for a complete activated sludge response without treating a different prediction task as a methodological failure.

\begin{table}[htbp]
\centering\footnotesize
\caption{Treatment prediction studies and the scope of their evidence.}
\label{tab:rrl_prediction_evidence}
\begin{tabular}{>{\raggedright\arraybackslash}p{0.24\linewidth}>{\raggedright\arraybackslash}p{0.33\linewidth}>{\raggedright\arraybackslash}p{0.31\linewidth}}
\toprule Study and task & Contribution established & Additional question for this dissertation\\\midrule
\citet{Ching2022}\newline Oxygen-demand soft sensing & Estimation across two plants with support for missing sensor observations. & Can a predictor estimate the complete effluent state from influent and proposed operating inputs?\\
\citet{Wang2025}\newline Effluent TN prediction & Comparison of six learned models with a mechanistic model and inspection of influential inputs. & Are the individual nitrogen and other component concentrations accurate and physically consistent?\\
\citet{Wang2024}\newline Influent component estimation & Neural estimation of five COD-related quantities used in process modeling. & How accurately is the resulting effluent component state predicted after treatment?\\
\citet{Fang2011}\newline Simulation-based operating search & Mechanistic simulation combined with a support vector surrogate and optimization. & Can connected component predictions be constrained and checked at the selected controls?\\
\citet{Selerio2026}\newline Complete reactor prediction & Structured prediction of 20 components with physical checks and training-size assessment. & How do common projection and plant connections affect prediction and operating assessment?\\
\bottomrule
\end{tabular}
\end{table}

\subsection{Model comparison, data availability, and extrapolation}

The methods used in treatment prediction provide several ways to approximate nonlinear responses. \citet{Breiman2001} developed random forests by averaging randomized trees. \citet{ChenGuestrin2016} developed XGBoost, while \citet{Ke2017} introduced the light gradient boosting machine (LightGBM). \citet{SmolaScholkopf2004} described support vector regression, which uses regularized fitting and kernel-based similarity. These methodological studies explain the learners used in wastewater applications. Their contributions support a comparator set with different fitting structures, while treatment studies establish which structures perform well for a particular process target.

Inspectable linear and nonlinear models provide another part of that comparison. \citet{Wold2001} described partial least squares regression, which uses latent components to summarize related input and output variation. These statistical components are weighted combinations of variables rather than physical wastewater constituents. \citet{Woo2009} applied a kernel extension to estimate COD, TN, and cyanide concentrations in industrial cokes wastewater treatment. It captured nonlinear process relationships and outperformed the assessed linear and other nonlinear partial least squares formulations. This establishes the benefit of the extension for that application. Ordinary partial least squares remains a different comparator for the present component target.

The tabular-model comparison of \citet{ShwartzZiv2022} and the neural-model survey of \citet{Borisov2024} provide further reasons to compare learning families rather than select one solely from its complexity. This is relevant to activated sludge data, where a row contains operating values and influent concentrations. However, the applicable model ranking must come from the shared treatment target and assessment design. The broader methodological literature guides comparator selection, while the treatment benchmark supplies process-specific evidence.

Training-size assessment addresses the cost of obtaining that evidence. \citet{VieringLoog2023} reviewed learning curves and found that their shapes depend on the learner, data, and task. A learning curve relates prediction performance to the number of training observations. \citet{Selerio2026} used such an assessment for its reactor surrogates. This addresses data efficiency within its stated input range. A broader comparison still needs to determine how the complete component errors and physical adjustments change together as simulated training cases become scarce.

Assessment also depends on information separation. \citet{Kaufman2012} defined information leakage and described how information unavailable for the intended prediction can enter model development. Applied to the treatment tasks reviewed above, this means that current effluent measurements can support a monitoring estimate but cannot automatically serve as inputs to a prediction made before treatment is evaluated. Model tuning and scaling also need to exclude the observations reserved for assessment. These are conditions for interpreting an error comparison. They do not imply that every study using downstream measurements has handled its data incorrectly.

The remaining comparative question combines model choice, training support, and operating domain. The reviewed indicator and component studies demonstrate useful prediction within their assessed conditions. A common complete-state comparison is needed to assess several learners under the same input and output definitions. It must also distinguish prediction within the training ranges from extrapolation beyond them. This motivates evaluating component error and physical validity together as influent and operating conditions change.

\section{Methods for imposing physical constraints on predictions}

\subsection{Physical guidance during fitting}

Physical constraints have been introduced into learning to address predictions that disagree with governing equations. \citet{Karniadakis2021} reviewed approaches that combine physical knowledge with machine learning. One influential approach is the residual-based fitting of \citet{Raissi2019}, which adds discrepancies in differential equations to the training objective. This makes equation disagreement costly alongside prediction error. The method provides a way to use governing equations when response observations are limited.

\citet{Li2024} examined this approach for water-treatment unit dynamics, including an activated sludge reactor. They compared physics-informed neural models with conventional learned models under limited data and different loading conditions. The physically informed models showed stronger predictive capability and generalization in those assessed systems. Their study therefore provides direct wastewater evidence that process equations can improve learning when data are scarce.

The distinction between this contribution and output enforcement follows from the fitting problem. A finite residual penalty balances equation agreement with other terms in the objective. It does not require every new predicted concentration vector to satisfy a balance equality exactly. Thus, the demonstrated learning benefit of \citet{Li2024} supports physical guidance during fitting, while sample-specific mass conservation and non-negativity require an additional construction or check. These are complementary requirements in activated sludge modeling.

\subsection{Conservation through model construction}

Other studies address conservation by changing how outputs are calculated. \citet{Beucler2021} introduced neural architectures that encode analytic constraints and compared them with constraints imposed through a loss function. In their climate-convection application, the architectural approach enforced the selected conservation laws to machine precision without degrading overall predictive performance. The study establishes that appropriately constructed outputs can satisfy known equalities directly.

\citet{SturmWexler2020} developed a conservation framework in which learned transfers are mapped to changes in chemical species through balanced process relationships. \citet{SturmWexler2022} then placed the stoichiometric mapping in a fixed neural layer for a photochemical surrogate. The later construction avoided the need to reconstruct all reaction transfers as separate training targets and retained atom conservation in the output changes. These studies address the coupling of species that independent output regression leaves unspecified.

Their relevance to activated sludge modeling lies in the shared use of stoichiometry. A balanced mapping can preserve the component combinations implied by the adopted reaction network. Concentration bounds remain a separate condition because a balanced change can consume more of a component than is present. This follows from the difference between a reaction-induced change and the final concentration. The construction must be assessed for the quantity it constrains.

\citet{Harder2022} explored both physical penalties and constrained output completion in an aerosol microphysics surrogate. Their assessment treated mass conservation and positive mass as explicit modeling concerns. It illustrates why conservation and concentration bounds need to be examined together when outputs are completed or adjusted. For the activated sludge task, the required state must satisfy both conditions after the complete prediction procedure has been applied.

\subsection{Projection after statistical prediction}

Projection provides an alternative that can be applied after a model has produced its output. \citet{SturmSilva2024} developed a species-weighted least-squares correction for atom conservation in an atmospheric chemistry surrogate. The weights accounted for the magnitude and uncertainty of individual species changes. In their assessment, unweighted correction could damage prediction of a small radical, while weighting retained its accuracy and slightly improved overall agreement. Their results establish that the projection metric affects prediction quality even when the enforced conservation relation is unchanged.

The corresponding activated sludge issue is variation in concentration scale across components. A numerical adjustment that is small for abundant substrate can be large relative to a residual nutrient concentration. The study of \citet{SturmSilva2024} therefore supports examining adjustment weights and component errors together. Its equality correction addresses atom accounting, while a simultaneous concentration-bound construction is needed when non-negativity is also required.

\citet{KircherVotsmeier2025} addressed that joint requirement through a positivity-preserving projection and linear interpolation backtracking. They evaluated atmospheric chemistry, heterogeneous catalysis, and synthetic reaction networks. The corrected predictions satisfied atom balances and remained positive in the tested systems without reducing overall prediction accuracy. This study establishes a method for imposing both requirements across different underlying surrogates. The remaining activated sludge question concerns its behavior under wastewater component units, reactor boundaries, training scarcity, and extrapolated operating inputs.

\citet{Selerio2026} supplied a direct activated sludge example. Its structured regression used physical penalties during fitting, followed by a check of each raw prediction and projection onto the mass conservation equalities. A linear program, which optimizes a linear objective under linear constraints, was used when non-negativity still failed. The published assessment reported no final mass conservation or non-negativity violations. The raw regression did not satisfy both requirements, so the final enforcement procedure was responsible for the physical result. This resolves the question of whether a structured single-reactor surrogate can return a component state with both imposed properties.

The same study also reported an accuracy trade-off relative to some unconstrained comparators. Consequently, its result establishes physical reliability for the returned state without making that state the most accurate prediction under every reported score. The next comparative question is how a common projection changes different learners when their raw predictions are held fixed. A paired comparison of raw and projected states can identify that effect across components and operating conditions.

Figure~\ref{fig:rrl_enforcement_locations} brings together the approaches described above. Fitting penalties, constrained output construction, and final-state projection address physical consistency at different stages. The distinction follows from the methods of the cited studies and explains which part of a predictor is responsible for a stated constraint.

\begin{figure}[htbp]
\centering
\begingroup\linespread{1}\selectfont
\input{figures/ch2_concept_enforcement_locations.tex}
\endgroup
\caption{Stages at which the reviewed methods introduce physical information. Fitting penalties guide equation agreement. Output construction and projection impose the relations encoded in their formulations. Simultaneous mass conservation and non-negativity require an assessment of the final component state.}
\label{fig:rrl_enforcement_locations}
\end{figure}

\subsection{Constraint scope and recent extensions}

More recent work connects physical enforcement to uncertainty and model capacity. \citet{Hansen2024} used integral conservation relations to adjust a probabilistic prediction, accounting for both predicted means and uncertainty. Their framework enforced conservation exactly when the imposed balance was treated as having zero uncertainty and was evaluated on a family of transport equations. \citet{Utkarsh2025} trained a probabilistic predictor together with a differentiable projection and evaluated it on differential-equation and time-series problems. A differentiable calculation provides derivatives that can be used during fitting. These studies address uncertainty as part of constrained prediction, rather than only correcting a single concentration estimate.

\citet{Min2024} developed neural output layers for input-dependent linear and convex constraints and established approximation results under the stated mathematical conditions. A convex feasible set contains the straight line joining any two feasible states. Their work addresses the concern that hard constraints, which require outputs to satisfy the imposed conditions, may restrict model capacity. \citet{Valente2025} examined projection onto physical manifolds, which are sets of outputs satisfying selected physical laws. Their spring-mass and reactive-plasma examples extend output correction to nonlinear physical relations. These contributions broaden the available constraint methods. Their evaluated problems and mathematical conditions must still be matched to a wastewater application.

Taken together, the physical-learning literature has already addressed conservation through several constructions and joint positivity through specific methods. The gap for this dissertation is therefore not the absence of physical enforcement. It is the need to assess its accuracy, adjustment size, and computational cost for a common activated sludge component task and then extend the stated accounting conditions to connected plant outputs. The adopted constraints remain narrower than the full biological and settling equations.

\section{Interpretable modeling of treatment responses}

\subsection{Inspection of inputs and fitted relationships}

Interpretability research addresses the need to examine why a model predicts a particular response. \citet{Rudin2019} argued for using directly inspectable models in consequential applications when such models are available. In wastewater treatment, this need appears when an engineer must understand how influent composition or an operating setting relates to predicted quality. The argument supports inspection of the model used for the decision. It does not establish the biological meaning of every fitted coefficient.

The treatment studies reviewed earlier use different forms of inspection. \citet{Wang2025} combined TN prediction with feature-attribution analysis to identify important inputs. \citet{Ekinci2023} used feature selection and attribution in sludge prediction. These approaches address which observed variables contribute to a specified statistical output. They can guide process investigation and reduce unnecessary inputs. Their explanations concern the fitted indicator or sludge model rather than all component balances of a mechanistic treatment system.

Explicit equations offer a complementary approach. \citet{Nazlabadi2021} reviewed response-surface methods in biological wastewater treatment, where polynomial terms describe operating effects and interactions over the assessed range. \citet{Brunton2016} developed sparse equation discovery from candidate functions for nonlinear dynamical systems. The latter is a methodological precedent for inspecting named terms, while response-surface applications provide the treatment context. Together, they support visible algebraic relationships as a useful form of model inspection.

For activated sludge analysis, an aeration term, an influent substrate term, and their product have different interpretations. The product permits the fitted aeration response to depend on the substrate condition. Such terms help describe a treatment response within the fitted domain. Establishing a causal biological parameter would require a different identification problem and evidence about the underlying mechanism. This distinction is particularly relevant when inputs are correlated.

\subsection{Structured regression and the final projected state}

\citet{Selerio2026} addressed the need for an inspectable complete-component surrogate by separating operating effects, influent effects, curvature, interactions, and learned dependence among effluent components. Its explicit second-order terms allow the fitted input relationships to be examined in physical coordinates. Component coupling then contributes to the complete raw prediction, and a fixed composition map produces the reported water quality quantities. The study therefore provides a direct treatment model that combines visible statistical structure with final physical enforcement.

The publication also identifies limits on that interpretation. Its displayed coefficient relationships describe the raw regression before the final constraint correction. The mass conservation projection and linear-program stages can change the returned state. A concentration bound becomes active when a predicted component reaches the permitted boundary, and that change can affect local sensitivity. The study explicitly distinguishes its coefficient displays from derivatives of the final projected output. It also recognizes that different fitted coefficient combinations can produce similar predictions and identifies coefficient stability across resampled data as a further diagnostic.

\citet{Shen2023} reviewed differentiable modeling that connects learned relationships with process equations. Their geoscience examples show how combined models can support interpretation while retaining physical knowledge, and they identify verification of the physical significance of outputs as a continuing need. This perspective supports examining the complete calculation in a treatment surrogate. The specific distinction between raw coefficients and final constrained sensitivities, however, is supplied directly by the activated sludge study of \citet{Selerio2026}.

The remaining interpretive question is thus an extension of an established contribution. An explicit raw equation already makes its terms visible. A fuller assessment must connect those terms to component coupling, reporting, and the constraints that determine the final effluent. It must also examine whether the inferred response patterns remain stable when the training cases change. This motivates interpretation of the complete prediction procedure while retaining the distinction between statistical association and biological mechanism.

\section{Surrogate modeling for operating decisions}

\subsection{Mechanistic and learned routes to operating optimization}

Activated sludge operating optimization predates the recent use of machine learning surrogates. \citet{Araujo2013} optimized a steady-state process model and used sensitivity analysis to select variables for economic control. Their analysis linked active effluent constraints and remaining operating degrees of freedom to the control structure. \citet{Nazif2023} calibrated a treatment model to a plant in Tehran, grouped influents by their quantitative and qualitative characteristics, and optimized operating settings for each group. These studies establish mechanistic routes that connect treatment constraints, influent conditions, and operating choices.

Learned models provide another route when repeated response evaluation is costly. The integrated approach of \citet{Fang2011} used mechanistic simulations to train a support vector surrogate before optimizing the nutrient-removal system. \citet{KusiakWei2013} fitted neural models for airflow, effluent carbonaceous biochemical oxygen demand, and TSS from industrial data. They then used a multiobjective particle-swarm search to examine dissolved-oxygen settings. These studies demonstrate operating searches supported by learned treatment relationships. Their targets determine which aspects of the operating decision are evaluated.

Dynamic control studies provide evidence at a different time scale. \citet{Bernardelli2020} developed a neuro-fuzzy predictive controller for aeration and field-tested it at a large municipal plant. They reported improved effluent quality and energy savings in that application. \citet{Croll2023} reviewed reinforcement learning for wastewater control and discussed the approaches and challenges involved. These studies establish that learning can participate in changing treatment operation, while the present dissertation concerns fixed-input steady-state operating assessment. Their dynamic results are relevant background, but they do not supply the evaluation of that steady-state task.

The operating literature therefore addresses both process-based and learned decision routes. The question for a new surrogate is how its selected settings compare on a common engineering basis. This requires the same treatment objectives, operating bounds, and reference calculations. It also requires distinguishing the cost of developing a learned model from the cost of using it during a search.

\subsection{Integrated resource recovery and connected responses}

Network and resource-recovery studies extend the decision beyond one biological unit. \citet{PadronPaez2020} examined sustainable treatment design through multiple objectives, while \citet{Aboagye2021} combined wastewater-network design, optimization, and sustainability assessment. Their contributions show how treatment and resource priorities can be evaluated together when several process choices are available. They establish a broader systems context for the operating objective of a connected plant.

\citet{Durkin2024} developed surrogate-based process synthesis for resource recovery from brewery wastewater. Their framework used high-fidelity simulations, regression surrogates, and classification models for process feasibility. The application considered biogas recovery, recoverable nutrient quality, and aeration demand, together with uncertainty in wastewater composition. It demonstrates that surrogates can support integrated decisions about several resource pathways while accounting for simulation feasibility.

This work addresses a network-design problem with alternative process configurations. The connected operating task here retains a specified activated sludge topology and predicts the mixer, reactor stages, clarifier outlets, and stored solids jointly. The additional requirement is consistency of those predicted component quantities at fixed physical connections. A classification of simulator feasibility and an explicitly imposed component balance answer different questions. The contribution of \citet{Durkin2024} motivates integrated decision analysis, while the reactor and clarifier literature defines the balances needed for the present extension.

The single-reactor result of \citet{Selerio2026} supplies physical enforcement at one unit. The connected mechanistic benchmark of \citet{Alex2008} supplies a precedent for comparing whole-plant responses and controls. Combining these lines of work leads to a specific remaining question. Can a statistical plant response preserve declared unit connections and solids inventories while supporting useful operating choices? This requires assessment beyond single-reactor feasibility and beyond the accuracy of one operating objective.

\subsection{Approximation error at selected operating conditions}

Surrogate optimization research addresses the reliability of the response used during search. \citet{BhosekarIerapetritou2018} reviewed surrogate modeling for feasibility analysis and optimization. \citet{McBride2019} reviewed surrogate construction from experiments and simulations for process-system decisions. Both place model development within a larger optimization problem. Their relevance is that a surrogate must approximate the quantities and operating region used by the decision, not only reproduce an average response over its training design.

\citet{Alexandrov1998} developed a trust-region framework that manages approximation models during optimization. The method uses reference information to assess a local approximation and control the region in which it is used. This is a methodological response to inaccurate surrogate behavior during search. For treatment analysis, it provides a reason to check the approximation near the controls actually selected rather than rely entirely on a prediction score averaged over other cases.

\citet{Pedrozo2025} compared five surrogate approaches for an optimization problem involving pooling of carbon-dioxide streams. They assessed predictive accuracy, computational effort, and optimization behavior separately. Their comparison of one-shot optimization with a trust-region filter strategy showed that the models differed in convergence and reliability during search. This is process-systems evidence that a favorable fitting result and a favorable optimization result are distinct outcomes. Its evaluated system is carbon capture, so the corresponding activated sludge decision still requires a treatment-specific comparison.

Mass conservation and non-negativity exclude certain impossible component states, but several feasible states can have different nitrogen distributions or biomass concentrations. An optimizer can therefore select a physically admissible prediction that remains inaccurate. The physical-enforcement studies establish the relevant constraints, while approximation-management studies establish the need for checking the response used during search. Together, they motivate independent mechanistic evaluation of the selected aeration, recycle, and wasting settings.

The remaining decision gap concerns the combined evidence for a connected activated sludge surrogate. Its unit balances and stored solids need to be consistent, its selected settings need to meet the original treatment requirements, and its objective needs to be evaluated on the same mechanistic basis as the direct comparator. Search time and verification effort must also be reported on a defined basis. This makes decision quality and computational benefit testable without assuming that constraint satisfaction alone establishes either outcome.

\section{Synthesis of established contributions and remaining gaps}

\subsection{What the reviewed studies have addressed}

The literature establishes a clear progression. Mechanistic activated sludge models define components, reaction rates, separation, and recirculation. Calibration studies connect those quantities to plant observations. Soft sensors and indicator models improve estimation of selected treatment responses. Simulation-based surrogates extend learned prediction to operating calculations. Physical-learning studies then address equation guidance and explicit constraint enforcement. Interpretable regression and surrogate optimization connect these predictions to inspection and decision use.

Several requirements have already been addressed directly. \citet{KircherVotsmeier2025} established joint atom balance and positivity for specified reactive-chemistry surrogates. \citet{Selerio2026} established final mass conservation and non-negativity for a structured activated sludge reactor surrogate and identified the distinction between raw and projected interpretation. \citet{Durkin2024} demonstrated integrated wastewater resource-recovery optimization using surrogates. These achievements define the starting point for the dissertation rather than requirements that the literature has entirely overlooked.

The remaining gaps arise when these contributions are brought into the same complete-component activated sludge assessment. Table~\ref{tab:rrl_evidence_gap_matrix} identifies the studies that address each need and the more specific question that remains. The questions concern the reviewed evidence and the declared reactor and plant setting. They do not claim that the underlying methods have never been developed.

\begin{table}[htbp]
\centering\footnotesize
\caption{Established studies and the remaining questions addressed by the dissertation.}
\label{tab:rrl_evidence_gap_matrix}
\begin{tabular}{>{\raggedright\arraybackslash}p{0.22\linewidth}>{\raggedright\arraybackslash}p{0.34\linewidth}>{\raggedright\arraybackslash}p{0.32\linewidth}}
\toprule Research need & Studies addressing the need & Remaining question\\\midrule
Complete treatment prediction & \citet{Wang2025} assesses TN prediction. \citet{Wang2024} estimates influent components. \citet{Selerio2026} predicts complete reactor components. & How do several learners compare on one complete-state target under training scarcity and extrapolation?\\
Physical validity of predictions & \citet{Beucler2021} and \citet{SturmWexler2022} impose conservation. \citet{KircherVotsmeier2025} and \citet{Selerio2026} address conservation and concentration bounds together. & How does common enforcement affect component accuracy, adjustment size, and cost across learners?\\
Interpretation of treatment responses & \citet{Wang2025} inspects input contributions. \citet{Selerio2026} exposes regression terms and distinguishes raw from final interpretation. & How do coupling, reporting, and active constraints affect sensitivities and stability of the explanation?\\
Connected operating decisions & \citet{Fang2011} and \citet{Durkin2024} use surrogates in treatment decisions. \citet{Alex2008} defines a connected plant benchmark. & Can joint component enforcement support settings that remain credible under independent plant evaluation?\\
\bottomrule
\end{tabular}
\end{table}

\subsection{The four linked questions of the dissertation}

The first gap concerns comparative prediction of the complete reactor response. The indicator studies establish useful outputs, and the structured-regression study extends prediction to all modeled components. The remaining assessment requires a shared component target across learning families and a clear account of how errors change with training size and inputs beyond the sampled range. This extends established prediction work to a common comparison rather than treating an indicator score as evidence about every component.

The second gap concerns the consequences of physical enforcement. Joint mass conservation and concentration bounds have already been established by specific methods. What remains is a paired assessment of their effect on the same raw predictions across several activated sludge learners. Component error, residuals, negative concentrations, adjustment size, and cost answer different questions. Their joint evaluation identifies the value and limits of enforcement without assuming an accuracy improvement.

The third gap concerns interpretation of the returned treatment state. The structured reactor regression already supplies visible terms and identifies the limits of raw coefficient inspection. The further question is how operating and influent sensitivities pass through coupling, composite reporting, and projection, including changes when concentration bounds become active. Examining stability across training partitions adds evidence about the reproducibility of those fitted relationships.

The fourth gap concerns connected prediction and operating assessment. Existing mechanistic plant models establish the shared streams and solids storage, while surrogate applications establish learned operating and resource-recovery decisions. The remaining requirement here is a joint constrained component response for the specified plant and an independent calculation of treatment outcomes at its selected controls. Comparison with direct mechanistic optimization then assesses the decision and computational cost under common priorities.

Figure~\ref{fig:rrl_intersecting_gaps} shows why these questions are linked. Physical enforcement acts on the concentrations being scored and interpreted. Those same concentrations enter the operating objective and constraints. A change in projection can therefore affect component errors, local sensitivities, and selected settings. The integrated assessment must follow that shared response through each stage.

\begin{figure}[htbp]
\centering
\begingroup\linespread{1}\selectfont
\input{figures/ch2_concept_intersecting_gaps.tex}
\endgroup
\caption{Links among the remaining comparative, physical, interpretive, and operating questions. Each question concerns the same defined activated sludge response. Established methods address parts of the problem, while their combined use requires common component and decision assessment.}
\label{fig:rrl_intersecting_gaps}
\end{figure}

The literature therefore supports the dissertation's progression from a defined reactor response to connected operating decisions. Chapters~\ref{ch:foundations} and \ref{ch:projection} develop the common component foundation and paired prediction assessment. Chapter~\ref{ch:interpretable} develops the structured regression and follows interpretation through physical enforcement. Chapter~\ref{ch:plant} extends the response to connected units and independently assesses selected controls. Chapter~\ref{ch:assessment} brings the evidence together within the stated simulation and modeling limits.
